\chapter{Resultados y discusión}
\label{chap:resultados}

\ph{Mientras el Capítulo~\ref{chap:pruebas} verificaba que el sistema está
bien construido, este capítulo evalúa si resuelve el problema planteado en el
Capítulo~\ref{chap:intro} y en qué medida: presenta los resultados obtenidos,
los contrasta con los criterios de éxito fijados, los compara con los
trabajos previos y los interpreta críticamente. Es el capítulo donde el
tribunal busca la evidencia de lo conseguido: cada afirmación debe apoyarse
en una tabla o una figura.}

\section{Diseño de la evaluación}
\label{sec:diseno_evaluacion}

\ph{Antes de mostrar números, define qué se mide y cómo. Retoma la
Tabla~\ref{tab:criterios-exito} del Capítulo~\ref{chap:intro}: sus criterios
de éxito son el guion de esta evaluación, y cada uno debe quedar respondido
en la Sección~\ref{sec:resultados_objetivos}. Describe:}

\begin{itemize}
  \item \textbf{\ph{Entorno de ejecución}}: \ph{hardware (CPU, RAM\dots) y
        software (sistema operativo, versiones relevantes) sobre los que se
        han obtenido las medidas, para que sean reproducibles.}
  \item \textbf{\ph{Datos o casos de estudio}}: \ph{origen, tamaño y
        representatividad de los datos, escenarios o usuarios empleados en la
        evaluación.}
  \item \textbf{\ph{Repeticiones}}: \ph{si hay medidas de rendimiento, indica
        cuántas veces se repite cada medición y por qué: una sola ejecución
        está sujeta a variabilidad (carga del sistema, caché, red\dots) y no
        es concluyente.}
  \item \textbf{\ph{Agregación}}: \ph{cómo se resumen los datos: media y
        desviación típica, mediana, percentiles (p.\,ej.\ p95 para tiempos de
        respuesta)\dots y por qué ese estadístico es el adecuado. Toda media
        debe ir acompañada de su desviación típica o de un intervalo de
        confianza; y si se comparan alternativas, conviene apoyar la
        comparación en un test estadístico sencillo (t-test o Wilcoxon).}
\end{itemize}

\ph{Rigor: distingue siempre lo que se ha \emph{medido} de lo que es una
\emph{impresión subjetiva}. «La interfaz es intuitiva» no es un resultado;
«8 de 10 usuarios completaron la tarea sin ayuda en menos de 2 minutos» sí
lo es.}

\section{Resultados}
\label{sec:resultados_objetivos}

\ph{Organiza esta sección con una subsección por objetivo (o por criterio de
éxito), en el mismo orden en que se enunciaron en el
Capítulo~\ref{chap:intro}. En cada subsección sigue el patrón: (1) presenta
el resultado en una tabla o figura, (2) referénciala desde el texto y (3)
interprétala frente al umbral fijado. Recuerda el principio: toda tabla o
figura se referencia desde el texto y se interpreta — un resultado sin
interpretación no es un resultado. Si el criterio de un objetivo ya quedó
verificado en un capítulo anterior (columna Cap.\ de la
Tabla~\ref{tab:criterios-exito}), su subsección se limita a recapitular esa
evidencia con \textbackslash ref y a dar el veredicto frente al criterio,
sin repetir datos. Las dos subsecciones siguientes modelan ambos patrones:
O1 ejemplifica la recapitulación y O2 la medición completa.}

\subsection{\ph{Resultados sobre el objetivo O1: título descriptivo}}
\label{subsec:resultados_o1}

\ph{Patrón RECAPITULACIÓN, para objetivos cuyo criterio ya quedó verificado
en un capítulo anterior según la columna Cap.\ de la
Tabla~\ref{tab:criterios-exito}. Siguiendo el ejemplo conductor, el criterio
de O1 (cobertura de todos los RF de prioridad alta) quedó verificado en los
Capítulos~\ref{chap:requisitos} y~\ref{chap:pruebas}: recapitula esa
evidencia con \textbackslash ref —la matriz de trazabilidad de la
Tabla~\ref{tab:trazabilidad_rf_pruebas} muestra que cada RF tiene al menos
una prueba que lo cubre, y la Tabla~\ref{tab:resumen_suite} acredita que la
suite completa pasa— y cierra con el veredicto frente al criterio («el
criterio de éxito asociado al objetivo O1 se considera cumplido»), sin
repetir datos ni añadir tablas nuevas.}

\subsection{\ph{Resultados sobre el objetivo O2: título descriptivo}}
\label{subsec:resultados_o2}

\ph{Patrón MEDICIÓN COMPLETA, para objetivos cuyo criterio se verifica en
este capítulo. Indica qué se ha medido para este objetivo y en qué
condiciones (remite a la Sección~\ref{sec:diseno_evaluacion}, no las
repitas). La Tabla~\ref{tab:resultados_o2} resume las métricas frente a sus
umbrales.}

\begin{table}[H]
\centering
\small
\begin{tabular}{@{}P{4.5cm} P{2.5cm} P{4.5cm} P{2cm}@{}}
\toprule
\textbf{Métrica} & \textbf{Valor} & \textbf{Umbral (RNF / criterio)} & \textbf{¿Cumple?} \\
\midrule
\ph{Tiempo medio de respuesta} & \ph{152~ms} & \ph{$\leq$ 200~ms (RNF-01)} & \ph{Sí} \\
\ph{Métrica 2} & \ph{\dots} & \ph{\dots (Tabla~\ref{tab:criterios-exito})} & \ph{\dots} \\
\ph{Métrica 3} & \ph{\dots} & \ph{\dots} & \ph{\dots} \\
\bottomrule
\end{tabular}
\caption{\ph{Resultados de las métricas asociadas al objetivo O2.}}
\label{tab:resultados_o2}
\end{table}

% Hueco para la gráfica del resultado. Descomenta y adapta:
% \begin{figure}[H]
%   \centering
%   \includegraphics[width=0.8\textwidth]{resultados_o2}
%   \caption{Pie autoexplicativo: qué se muestra, con qué datos y qué se observa.}
%   \label{fig:resultados_o2}
% \end{figure}

\ph{(Espacio reservado para una gráfica del resultado. Elige el tipo
adecuado al dato: líneas para evolución o escalabilidad, barras para comparar
alternativas, diagramas de caja para mostrar dispersión. Pie autoexplicativo
y referencia desde el texto.)}

\ph{Interpretación crítica contra el umbral, no valoración vaga: en lugar de
«el rendimiento salió bien», escribe «el valor de 152~ms queda por debajo del
umbral de 200~ms fijado en el RNF-01 (Tabla~\ref{tab:rnf}), por tanto el
criterio de éxito asociado al objetivo O2 se considera cumplido». Si una
métrica no alcanza el umbral, dilo igual de claro y analiza la causa:
reconocerlo vale más ante el tribunal que ocultarlo.}

\subsection{\ph{Resultados sobre el objetivo O3: título descriptivo}}
\label{subsec:resultados_o3}

\ph{Elige para cada objetivo el patrón que corresponda (recapitulación o
medición completa) y añade tantas subsecciones como objetivos tenga el
trabajo; ningún objetivo del Capítulo~\ref{chap:intro} debe quedarse sin
resultado.}

\section{Comparación con trabajos previos}
\label{sec:comparacion_trabajos}

\ph{Retoma la Tabla~\ref{tab:comparativa_relacionados} del
Capítulo~\ref{chap:estado_arte}, ahora con tus resultados en la mano: ¿qué
celdas marcadas como «Parcial» o «No» en los otros trabajos cubre ya tu
solución, con qué evidencia de la Sección~\ref{sec:resultados_objetivos}? ¿Y
en qué queda tu trabajo por detrás de ellos? Una comparación honesta en ambas
direcciones es lo que distingue un análisis riguroso de un alegato.}

\ph{Si los trabajos relacionados publican métricas comparables (mismos datos
o condiciones equiparables), añade aquí una tabla numérica con una fila por
métrica y una columna por trabajo, citando la fuente de cada
cifra~\cite{apellido2024articulo}. Las cifras propias deben ir acompañadas
de su dispersión (desviación típica o intervalo de confianza), no solo de la
media. Si no existen métricas comparables, no la
fuerces: justifica por qué la comparación directa no es posible (datos,
entorno o alcance distintos) y mantén la comparación cualitativa por
criterios.}

\section{Discusión}
\label{sec:discusion}

\ph{Interpretación global, sin repetir los números anteriores: (1) ¿qué
significan los resultados en conjunto? ¿responden al problema de la
Sección~\ref{sec:problema}? (2) ¿qué resultados fueron inesperados y a qué
los atribuyes (hipótesis razonadas, no excusas)? (3) ¿qué compromisos
(\emph{trade-offs}) revelan: rendimiento frente a simplicidad, coste frente a
precisión, generalidad frente a especialización\dots?}

\ph{Cierra con un párrafo breve de amenazas a la validez: qué podría
debilitar estas conclusiones, distinguiendo validez \emph{interna} (factores
del propio experimento: datos sintéticos, selección de casos, configuración),
\emph{externa} (hasta qué punto generalizan los resultados a otros entornos,
datos o usuarios) y \emph{de constructo} (¿miden las métricas elegidas
realmente lo que se quiere evaluar?). Basta un párrafo: el tratamiento
completo de las limitaciones se hace en la Sección~\ref{sec:limitaciones},
no lo dupliques aquí.}

% --- Transición al siguiente capítulo ---
\ph{Cierra el capítulo con un párrafo puente de 2--4 líneas: qué queda
establecido aquí y qué pregunta abre el capítulo siguiente. Por ejemplo:
«Evaluados los resultados frente a los criterios de éxito y discutido su
alcance, el Capítulo~\ref{chap:conclusiones} sintetiza el grado de
cumplimiento de los objetivos planteados en el Capítulo~\ref{chap:intro},
las limitaciones del trabajo y las líneas de trabajo futuro que estos
resultados abren».}
